\documentclass[11pt,letterpaper]{article}
\usepackage[margin=1in]{geometry}
\usepackage{amsmath,amssymb,amsthm}
\usepackage{physics}
\usepackage{hyperref}
\usepackage{booktabs}

\numberwithin{equation}{section}
\renewcommand{\theequation}{S18.\arabic{equation}}

\title{Supplement S18: Neutrino Mass from Hodge Duality \\
on $\Lambda^{\bullet}(\mathbb{R}^{9})$ \\
\large Supporting Manuscript \S XIV}
\author{UDT Collaboration}
\date{}

\begin{document}
\maketitle

\begin{abstract}
The neutrino mass formula $m_\nu = \alpha^3 m_e/4$ is derived from
the Hodge duality structure of the exterior algebra
$\Lambda^\bullet(\mathbb{R}^9)$, where $9 = 2\ell + 2|\kappa_{\max}|
+ 1$ is the total angular dimension from the Diophantine triple.
Each factor of $\alpha$ corresponds to one electromagnetic vertex
(grade 2 on $\Lambda^\bullet$), and three vertices are the minimum
required to reach grade~6 --- the Hodge dual of the pion (grade~3).
\end{abstract}

\section{The Exterior Algebra}

The Diophantine triple $(j, \ell, |\kappa_{\max}|) = (1/2, 1, 3)$
generates a total angular space of dimension
$n = 2\ell + 2|\kappa_{\max}| + 1 = 9$.  The exterior algebra
$\Lambda^\bullet(\mathbb{R}^9)$ has total dimension $2^9 = 512$, with
grade-$k$ subspaces of dimension $\binom{9}{k}$.

The physically significant grades:
\begin{center}
\begin{tabular}{ccll}
\toprule
Grade $k$ & $\dim\Lambda^k$ & Role & Particle \\
\midrule
0 & 1 & Scalar & Vacuum \\
2 & 36 & Antisymmetric coupling & $\alpha_{\rm EM}$ ($36 = 4 \times 9$) \\
3 & 84 & Symmetric occupation & Pion ($84\pi m_e$) \\
6 & 84 & Hodge dual of grade 3 & Neutrino \\
\bottomrule
\end{tabular}
\end{center}

\section{Hodge Duality and the Neutrino}

The Hodge star maps $\Lambda^k \to \Lambda^{9-k}$.  In particular:
\begin{equation}
  *\Lambda^3 = \Lambda^6, \qquad \dim\Lambda^3 = \dim\Lambda^6 = 84\,.
\end{equation}
The pion occupies grade~3 (the bosonic occupation count
$\binom{9}{3} = 84$).  Its Hodge dual at grade~6 is the neutrino
sector.

\section{The Three-Vertex Mechanism}

Each electromagnetic coupling corresponds to grade~2 on
$\Lambda^\bullet$.  The wedge product
$\Lambda^2 \wedge \Lambda^2 \wedge \Lambda^2 \to \Lambda^6$ maps
three EM vertices to the neutrino grade.  This is topological:
\emph{one cannot reach grade~6 with fewer than 3 copies of
$\Lambda^2$} (since $2 \times 2 = 4 < 6$, two vertices only reach
grade~4).

Each vertex contributes one factor of $\alpha_{\rm EM}$:
\begin{equation}
  m_\nu = \alpha_{\rm EM}^3 \times m_e \times
  \frac{\text{source}(\kappa\!=\!-1)}{1}
  = \frac{\alpha^3 m_e}{4}\,,
\end{equation}
where source$(\kappa\!=\!-1) = 1/4$ is the ground-state Dirac source
integral (the orbit-matching factor).

\section{Cross-Checks}

Three independent algebraic identities verify the result:
\begin{align}
  \frac{m_p \times m_\nu}{m_e^2}
  &= \frac{3}{2}\pi^5 \alpha^3 &&\text{(exact)}\,,\\
  \frac{m_\mu}{m_\nu}
  &= \frac{80}{3}\left(\frac{\pi}{\alpha}\right)^{\!3}
  &&\text{(exact)}\,,\\
  \alpha^3 m_e \times S(-1)
  &= m_\nu &&\text{(source connection)}\,.
\end{align}
All three hold to machine precision, confirming that $m_\nu$ is not a
numerical coincidence but a structural consequence of Hodge duality on
the angular exterior algebra.

\section{Prediction}

\begin{equation}
  m_\nu = \frac{\alpha^3 m_e}{4} = 0.049\;\text{eV}
  \qquad(\sqrt{\Delta m^2_{31}} = 0.050\;\text{eV},\ {-0.8\%})\,.
\end{equation}
Zero free parameters.  Testable by KATRIN and Project~8.

\end{document}
